\section{Hypothesis register, amendments and deviations}
\label{app:register}

The analysis plan was fixed in version control before the analyses were run, with ten hypotheses (H1--H10)
and their thresholds. It was amended \resAmendmentCount{} times. This appendix gives every hypothesis with its
rule, its origin and the status computed from the stored results, and then each amendment with its reason.
No numerical threshold of any earlier document was changed by any amendment, and no earlier verdict was edited;
later analyses add rows.

{\small
\begin{longtable}{llp{0.55\linewidth}l}
\toprule
& origin & rule & status \\
\midrule
\endhead
H1 & plan & Report the tool output verbatim; if k\_e and k\_a are globally identifiable, delete all swap framing. & formulation-dependent \\
H2 & plan & Ratio $\le$ 0.02 (linearized, 6 h window); median ratio $\le$ 0.25 (nonlinear, 3 h window). & met \\
H3 & plan & Median Schur-complement condition number decreases iAUC $\to$ iAUC+centroid $\to$ trace, paired Wilcoxon per step, Holm, alpha 0.05 (1x box, each objective's own estimate). & not met \\
H4 & plan & Median |cos| of the weak timing eigenvector with (1/k\_a, -1/k\_e), iAUC+centroid, $\ge$ 0.8. & met \\
H5 & plan & k\_e and k\_a bounded in $\le$ 25\% of subjects under iAUC and $\ge$ 50\% under the trace; S\_I bounded in $\ge$ 80\% under iAUC among subjects interior for S\_I (1x box). & not met \\
H6 & plan & Median per-subject Spearman between the gradient diagnostic and the profile-likelihood ranking $\ge$ 0.7; AUC $\ge$ 0.85 for detecting profile-flat parameters in the synthetic study. & partially evaluated \\
H7 & plan & Paired difference in held-out iAUC MAE, gradient vs exhaustive grid (three parameters and S\_I only), 90\% CI inside (-150, +150) mg/dL*min. & met \\
H8 & plan & Gradient three parameters vs S\_I only: 90\% CI inside (-150, +150). & met \\
H9 & plan & A full-trace fit lowers held-out iAUC MAE relative to an iAUC fit by more than 150 mg/dL*min (paired Wilcoxon, Holm). & not met \\
H10 & plan & Kendall tau $\ge$ 0.8 across 5 inits, 3 bound settings, parameterizations, optimizers; S\_I ranked first in $\ge$ 95\% of interior subjects. & not evaluated \\
H11 & amendment one & tau1 CI bounded in $\ge$ 70\% (iAUC+centroid and trace); p bounded $\le$ 30\% (iAUC+centroid), $\ge$ 50\% (trace); among subjects interior for the parameter. & not met \\
H12 & amendment one & Tied-rate model within +/-150 of the three-parameter model (90\% CI); tau1 bounded in $\ge$ 70\% under iAUC+centroid. & not met \\
H13 & amendment four (post hoc) & Replica with the model true: S\_I interval bounded in $\ge$ 50\% of interior subjects, k\_e and k\_a in $\le$ 25\%. & not met \\
H14 & amendment four (post hoc) & Replica held-out iAUC MAE / real MAE: $\le$ 0.7 substantial part unexplained by noise or optimization; $\ge$ 0.9 noise accounts for almost all; between, mixed. & reported (no pass/fail) \\
H15 & amendment four (post hoc) & Shanghai and Hall (iAUC, primary box): k\_e, k\_a bounded in $\le$ 25\%, weak Fisher direction cosine $\ge$ 0.8 with the exchange direction in $\ge$ 80\% of subjects. S\_I: no threshold. & met \\
H16 & amendment four (post hoc) & Exploratory: S\_I saturation curve; carbohydrate-scale sensitivity; heteroskedastic noise (no thresholds). & reported (exploratory) \\
H17 & amendment five (post hoc) & Replica seeds 0 to 4 pooled: the H13 rule (S\_I bounded in $\ge$ 50\% of interior subjects, k\_e and k\_a in $\le$ 25\%); reported with the per-seed range and the coverage of the bounded intervals. & not evaluated \\
H18 & amendment five (post hoc) & Coordinate profiles on the true-model replica (seed 0, polished estimates): under the trace tau1 bounded in $\ge$ 70\% and p in $\ge$ 50\% of interior subjects. & not evaluated \\
H19 & amendment five (post hoc) & Synthetic recovery with random truths (seeds 0 and 1 pooled): the H6 thresholds on the synthetic data (median Spearman $\ge$ 0.7, AUC $\ge$ 0.85). & not evaluated \\
H20 & amendment five (post hoc) & Shanghai, A9 protocol: gradient vs grid and three parameters vs S\_I only, 90\% CI inside (-150, +150) on held-out iAUC MAE. & not evaluated \\
H21 & amendment five (post hoc) & Coordinate fits: the Adam estimate is a likelihood optimum to within 1.92 in $\ge$ 90\% of subjects (iAUC+centroid and trace), against L-BFGS from five random starts. & not evaluated \\
H22 & amendment five (post hoc) & Dalla Man model, iAUC: kabs, kmax and kmin each bounded in $\le$ 25\% of subjects (profile intervals, primary box). & not evaluated \\
\bottomrule
\caption{The hypothesis register: rule as written in the plan or amendment, origin, and the status computed from the stored results.}
\label{tab:register}
\end{longtable}
}


\paragraph{Amendment one (before any identifiability analysis).}
Two secondary hypotheses were added (the identifiable combination $\tau_1$ and the tied-rate model, H11
and H12) and three method corrections adopted: every identifiability analysis uses the unregularized
maximum-likelihood estimate; the per-observable weights are inverse residual variances from a pilot fit; and
the trace objective is whitened for first-order autocorrelation. Reason: a penalty toward the population values,
weights taken from the spread of the observations, and independent-error treatment of autocorrelated sensor
residuals would each have made the parameters look better identified than the likelihood supports. The
multistart analysis was moved from the cross-validation to the full-data analysis of the inference gap, and
the grid cells were precomputed per meal; neither changes a hypothesis.

\paragraph{Amendment two (after the first Fisher results were viewed).}
The definition of an interior subject was changed from ``no parameter on a bound'' to a per-parameter
definition. Reason: parameters that are not identified sit on a bound in any box, so under the original
definition the set for the $\SI$ clause is empty by construction. The original definition is still
computed and reported, and the clause as written is recorded as unevaluable. The box used for the cross-method
comparison stays the one the baselines were trained on; the half-width, original and double-width boxes are
reported side by side without choosing one.

\paragraph{Amendment three (after the structural-identifiability tool had been run).}
The reporting of the structural result was fixed. It depends on the initial gut state given to the tool, so all
runs are reported and the verdict is labelled formulation-dependent. A sweep of initial gut content and a
check of the optimizer budget of the gradient cells were declared before any cross-validation result existed.

\paragraph{Amendment four (after the CGMacros results; post hoc).}
Hypotheses H13--H16 were declared after the $\SI$ clause of H5 had failed, in response to the question of
whether the weakness is inherent to the observable. They add the true-model replica (H13, H14), replication on
two cohorts (H15) and exploratory checks (H16).

\paragraph{Amendment five (after the first freeze of the results; post hoc).}
Hypotheses H17--H22 and the operational reading of H10 were declared after the first freeze of the results,
with their thresholds, before the corresponding analyses were run: replica stability over seeds (H17), the
combinations under a true model (H18), a synthetic study with random truths, which completes the synthetic
clause of H6 (H19), prediction on a second cohort (H20), whether the coordinate estimates are likelihood optima
(H21), and a second model class (H22). The amendment also records four corrections to the earlier record, none
of which changed a verdict: the sign of a descriptive field in the carbohydrate-scale analysis; a per-unit time
limit that left two units of one profile run unfinished (the run was repeated with a longer limit); the
stated cohort sizes, which are those the selection rule produces; and a summary section that read the most
complete result set without checking the cohort.

\paragraph{Where the plan and the record differ.}
The plan stated H5 as a rule over subjects ``interior-converged'', which is empty here (amendment two). The
plan's H1 rule fired under the planned formulation (amendment three). The error decomposition of an earlier
analysis, which had no definition under which its terms added, was not carried forward; the replica of H13--H14
replaces it with what-if noise budgets that make no additive claim.

\section{Analyses planned and not run}
\label{app:notrun}

The sensitivity of the $\SI$ profile to heteroskedastic observation noise, and the design sweep over meal size
and timing, were planned and not run. An analysis of clinical correlations of $\SI$ across cohorts was
not repeated; only the Shanghai check of Section~\ref{sec:leakage} is reported. \todo{List here any analysis of amendment five that did not finish before the second freeze.}
